%===============================================================
% 1. Определение \cabinetlist (все варианты)
%===============================================================
\providecommand{\cabinetlist}{}
\IfStrEqCase{\devicetype}{%
  {ЮНИТ-М500-ВЧ}{\renewcommand{\cabinetlist}{ШЭТ 220.01-0, ШЭТ 220.02-0, ШЭТ 220.03-0, ШЭТ 220.04-0, ШЭТ 221.01-1, ШЭТ 221.02-1 }}%
  {ЮНИТ-М500-ДЗЛ}{\renewcommand{\cabinetlist}{ШЭТ 220.05-0, ШЭТ 221.06-0, ШЭТ 221.03-1 }}%
}[\renewcommand{\cabinetlist}{\textcolor{red}{ШЭТ XXX.XX-X}}]

%===============================================================
% 2. Определение \functionlist (все варианты)
%===============================================================
\providecommand{\functionlist}{}
\IfStrEqCase{\devicetype}{%
  {ЮНИТ-М500-ВЧ}{\renewcommand{\functionlist}{Основные защиты линии (ДФЗ/НВЧЗ/ВЧБ), комплект ступенчатых защит (ДЗ, ТНЗНП, МТО, АМТЗ, ТЗО и др.), автоматика управления выключателем}}%
  {ЮНИТ-М500-ДЗЛ}{\renewcommand{\functionlist}{Основная защита линии (ДЗЛ), комплект ступенчатых защит (ДЗ, ТНЗНП, МТО, АМТЗ, ТЗО и др.), автоматика управления выключателем}}%
}[\renewcommand{\functionlist}{\textcolor{red}{Перечень защит}}]

%===============================================================
% 3. Раздел «Описание и работа»
%===============================================================
{\color{unidarkgreen}\section[Описание и работа]{Описание и работа}}

{\color{unidarkgreen}\subsection{Назначение}}

\begin{enumerate}
\item Устройство предназначено для установки в релейных залах электростанций и подстанций на панелях и в составе шкафов (в том числе \cabinetlist для установки на объектах филиалов и ДО ПАО «Россети»), построенных на базе архитектур I и II типа, а также на базе смешанной архитектуры.
\item Устройство разработано для применения в схемах вторичной коммутации распределительных устройств 110–220 кВ в составе комплексов РЗА, ПА, АСУ ТП, ССПИ на объектах энергетики.
Устройство предназначено для реализации функций релейной защиты и автоматики, при этом может осуществлять функции измерения, регистрации, осциллографирования и сигнализации.
При подключении устройства к информационной сети энергообъекта возможно осуществление функций дистанционного управления устройством и сбор данных с него посредством протоколов и интерфейсов, которые определяются конфигурацией устройства.
\item Аппаратная часть устройства определяется кодом заказа, представленным в приложении \ref{app:order}.
\item Взаимодействие с устройством может быть реализовано при помощи подключаемого блока интерфейсного «ЮНИТ-ИЧМ» или через ПК с установленным ПО «ЮНИТ-СЕРВИС-РЗА». Руководство по эксплуатации блока «ЮНИТ-ИЧМ» представлено в документе \ManualUnitHMI. Руководство пользователя представлено в документе \ManualUnitSRZA.
\item Подробное описание назначения устройства представлено в \ManualUnitGeneral.
\end{enumerate}

%===============================================================
% 4. Технические характеристики
%===============================================================
{\color{unidarkgreen}\subsection{Технические характеристики}}

\begin{enumerate}
\item Основные номинальные параметры устройства указаны в таблице \ref{desc:chars}.
\begingroup
\small
\setlength{\extrarowheight}{0.06cm}
\begin{longtable}{|>{\centering\arraybackslash}m{12.9cm}|>{\centering\arraybackslash}m{4cm}|}
\caption{Основные номинальные параметры устройства\hfill\vspace{-0.5\baselineskip}}\label{desc:chars}\\
\hhline{|-|-|}
\rowcolor{gray!30} Параметр & Значение \\ \hhline{|-|-|}
\endfirsthead
\caption*{\hspace{3pt}\emph{Продолжение таблицы \ref{desc:chars}\hfill\vspace{-0.5\baselineskip}}}\\
\hhline{|-|-|}
\rowcolor{gray!30} Параметр & Значение \\ \hhline{|-|-|}
\endhead
\hhline{|-|-|}
\endfoot
\endlastfoot
\raggedright Номинальный переменный ток, А & \centering\arraybackslash 1 или 5 \\ \hhline{|-|-|}
\raggedright Номинальное линейное напряжение переменного тока, В & \centering\arraybackslash 100 \\ \hhline{|-|-|}
\raggedright Номинальная частота, Гц & \centering\arraybackslash 50 \\ \hhline{|-|-|}
\raggedright Номинальное напряжение оперативного постоянного / переменного тока, В & \centering\arraybackslash 220 \\ \hhline{|-|-|}
\end{longtable}
\endgroup

\item Подробное описание основных технических характеристик конструктивного исполнения, диэлектрических свойств, электропитания устройства, входов измерительных цепей тока и напряжения, дискретных входов, входа 1PPS, выходных сигнальных реле, а также данные по электромагнитной совместимости, помехоэмиссии и надежности устройства представлены в \ManualUnitGeneral.
\end{enumerate}

%===============================================================
% 5. Конструктивное исполнение
%===============================================================
{\color{unidarkgreen}\subsection{Конструктивное исполнение}}

\begin{enumerate}
\item Устройство конструктивно имеет блочно-модульное строение и является неразборным (не требующим демонтажа блоков для выполнения обслуживания в процессе эксплуатации).
Конструкция устройства предусматривает навесной способ установки на монтажную плиту и обеспечивает возможность одностороннего обслуживания.
Устройство предусматривается к исполнению в трёх вариантах для различных архитектур:
\begin{itemize}
\item М510 (компактный) – 1/2 от 19 дюймов;
\item М514 (средний) – 3/4 от 19 дюймов;
\item М519 (максимальный) – 19 дюймов.
\end{itemize}
\item В устройстве устанавливается измерительный блок аналоговых сигналов. Типовая схема подключения входных аналоговых сигналов к устройству приведена в приложении \ref{app:analogs}. Представление выполнено для измерительного блока аналоговых сигналов «М179» в соответствии с кодом заказа для типового исполнения, в котором предусмотрено 7 токовых измерительных каналов с номинальным током 5 А и 9 измерительных каналов для измерения напряжения (в устройстве «ЮНИТ-М510» возможно размещение в слотах «М4», «М5»; в устройстве «ЮНИТ-М514» -- в слотах «М6», «М7», «М8» ; в устройстве «ЮНИТ-М519» -- в слотах «М10», «М11», «М12»).
\item Подробное описание конструкции, представление внешнего вида исполнений устройства, описание параметров настроек и технических характеристик блоков дискретных входов, блоков дискретного управления и измерительных блоков, а также габаритные и установочные размеры устройства приведены в документации \ManualUnitGeneral.
\end{enumerate}

%===============================================================
% 6. Функциональный состав
%===============================================================
{\color{unidarkgreen}\subsection{Функциональный состав}}

\begin{enumerate}[label=\arabic{section}.\arabic{subsection}.\arabic*, labelsep=7pt, leftmargin=0pt, itemindent=60pt, align=left]

\item Перечень функций, реализованных в составе устройства, приведен в таблице \ref{desc:tbl1}.
Описание функций релейной защиты и автоматики в составе устройства совместно с их функциональными схемами приводятся в разделе \ref{sec:funcs} данной документации.
Подробное описание вспомогательных (сервисных) функций приведено в документации \ManualUnitGeneral.
Общая функционально-структурная схема взаимосвязей функций приведена в приложении \ref{app:fsu}.
Перечень уставочных параметров функций в составе устройства представлен в приложении \ref{app:settings}.

\begingroup
\small
\setlength{\extrarowheight}{0.05cm}
\begin{longtable}{|>{\centering\arraybackslash}m{3.2cm}|>{\raggedright\arraybackslash}m{13.7cm}|}
\caption{Функции в составе устройства «ЮНИТ-М500-ВЧ»\hfill\vspace{-0.5\baselineskip}}\label{desc:tbl1}\\
\hhline{|-|-|}
\rowcolor{gray!30} Группа функций & Примечание \\ \hhline{|-|-|}
\endfirsthead
\caption*{\hspace{3pt}\emph{Продолжение таблицы \ref{desc:tbl1}\hfill\vspace{-0.5\baselineskip}}}\\
\hhline{|-|-|}
\rowcolor{gray!30} Группа функций & Примечание \\ \hhline{|-|-|}
\endhead
\hhline{|-|-|}
\endfoot
\endlastfoot
Основные функции & \functionlist \\ \hline
Сервисные функции &
  Сигнализация (обеспечение ввода служебных дискретных сигналов о состоянии оборудования шкафа РЗА, формирование отчетов по изменению состояния дискретных внешних и внутренних сигналов и выдачу отчетов в АСУ~ТП, формирование и выдачу сигналов аварийной и предупредительной сигнализации в схему резервной сигнализации)\newline
  Осциллографирование и журналирование (запись доаварийных, аварийных и послеаварийных режимов с сохранением осциллограмм в энергонезависимой памяти устройства, запись о событиях по изменению конфигурации устройства, изменения состояния внутренних сигналов алгоритмов, обновление встраиваемого программного обеспечения и пр. с метками времени и описанием события в энергонезависимой памяти устройства)\newline
  Самодиагностика (непрерывный оперативный контроль работоспособности в течение всего времени работ)\newline
  Конфигурирование (задание/переключение уставок, ввод в работу / вывод из работы, изменение параметров самоописания устройства, перезагрузка устройства, ввод/вывод тестового режима, изменение режима работы информационных сервисов, возможность переопределения связей входных и выходных логических сигналов алгоритмов устройства с дискретными входами, выходными реле, светодиодными индикаторами и функциональными клавишами управления)\newline
  Тестовый режим (обработка принятых сигналов и отправка отчетов с установленным флагом «Тест»)\newline
  Самоописание (предоставление данных об устройстве посредством устройства «ЮНИТ-ИЧМ» и сервисное программное обеспечение «ЮНИТ-СЕРВИС-РЗА»)\newline
  Синхронизация времени (календарь и часы реального времени, синхронизируемые от персонального компьютера или от АСУ~ТП) \\ \hhline{|-|-|}
\end{longtable}
\endgroup

\item Устройство поддерживает два режима оперативного управления: местный и дистанционный.
Выбор режима управления устройством производится посредством функциональной клавиши «М/Д» или по состоянию выбранного дискретного входа.
В местном режиме игнорируются все команды управления с верхнего уровня АСУ ТП. Более подробно режимы устройства описаны в документации \ManualUnitGeneral.

\item Для управления режимами работы основных функций (оперативный ввод/вывод, перевод на сигнал и пр.), а также для генерации различных сигналов сброса функций, применяются виртуальные ключи и кнопки.
Подробное описание и сценарии работы виртуальных ключей и кнопок приведено в документации \ManualUnitGeneral.

%-----------------------------------------------
% Группы уставок
%-----------------------------------------------
\item Группы уставок
\begin{enumerate}
\item Устройство имеет четыре группы уставок, предусмотрено оперативное переключение между несколькими сохраненными группами уставок. Подробное описание работы с группами уставок приведено в документации \ManualUnitGeneral.
\item Изменение группы уставок возможно:
\begin{itemize}
  \item по месту, через устройство «ЮНИТ-ИЧМ», п.\ref{sec1:hmi};
  \item дистанционно, через команды от АСУ ТП;
  \item дистанционно, средствами сервисного ПО «ЮНИТ-СЕРВИС-РЗА»;
  \item по месту, через соответствующий дискретный вход сигналом импульсного типа.
\end{itemize}
\item Переключение группы уставок происходит не мгновенно.
После появления сигнала/команды на изменение активной группы уставок и до момента активации новой группы проходит от 7 до 15 с.
В это время происходит проверка целостности файлов выбранной группы уставок.
Всё это время устройство продолжает функционировать на той группе уставок, которая была активна до появления сигнала/команды на изменение активной группы уставок.
\end{enumerate}

%-----------------------------------------------
% Аварийный осциллограф
%-----------------------------------------------
\item Аварийный осциллограф
\begin{enumerate}
\item При возникновении нарушений нормальных режимов работы данный процесс может записываться в память с помощью аварийного осциллографа.
Осциллографирование выполняется с частотой дискретизации не менее 1200 Гц.
Осциллограммы содержат измеренные значения аналоговых сигналов, дискретные сигналы состояния входов, выходных реле, а также логических сигналов функций РЗА и сервисные сигналы устройства.
Подробное описание параметров встроенного регистратора аварийных событий приведено в документации \ManualUnitGeneral.
\item Пуск осциллографа производится по факту изменения состояния дискретных сигналов, заданных в качестве критериев пуска.
В зависимости от длительности записи предаварийного и послеаварийного режима, в памяти может храниться разное количество осциллограмм.
Данные, собранные аварийным осциллографом, сохраняются в энергонезависимой памяти устройства и циклически перезаписываются.
Удаление записей аварийного осциллографа пользователем невозможно.
\item Перечень сигналов, доступных для регистрации аварийным осциллографом, приведен в приложении \ref{app:signals}.
\end{enumerate}

%-----------------------------------------------
% Журнал событий
%-----------------------------------------------
\item Журнал событий
\begin{enumerate}
\item В устройстве реализована функция записи сигналов дискретных входов, работы защит, автоматики и сервисных сигналов в журнал событий.
События безопасности записываются в отдельный журнал безопасности.
Подробное описание параметров журналов событий приведено в документации \ManualUnitGeneral.
\item Максимальное число записей журнала -- 20000.
Каждая запись в журнал событий содержит в себе следующую информацию:
\begin{itemize}
\item дату и время возникновения события;
\item идентификатор события;
\item состояние дискретного или значение аналогового сигнала в момент события;
\item качество зарегистрированного параметра;
\item качество времени.
\end{itemize}
\item Данные о событиях сохраняются в энергонезависимой памяти.
При заполнении журнала событий перезапись осуществляется от более старых событий к новым, то есть циклически.
Удаление информации из журнала событий пользователем невозможно.
Перечень сигналов, доступных для регистрации в журнале событий, приведен в приложении \ref{app:signals}.
\end{enumerate}

%-----------------------------------------------
% Режим тестирования
%-----------------------------------------------
\item Режим тестирования
\begin{enumerate}
\item В устройстве реализована поддержка режима тестирования «Тест» в объеме требований стандартов МЭК 61850.
Режим тестирования предусмотрен для возможности оперативной проверки элементов/функций, таких как: ФСУ, измерительные органы, дискретные входы и выходные реле устройства.
\item Особенности перевода устройства в режим тестирования и функциональности устройства в этом режиме описаны в документации \ManualUnitGeneral.
\end{enumerate}

%-----------------------------------------------
% Синхронизация времени
%-----------------------------------------------
\item Синхронизация времени
\begin{enumerate}
\item Устройство предусматривает возможность синхронизации внутренних часов реального времени несколькими способами:
\begin{itemize}
\item по протоколу SNTP;
\item по протоколу PTP;
\item через дискретный вход устройства с использованием сигнала синхронизации 1PPS;
\item средствами стандартизованного протокола передачи данных.
\end{itemize}
Подробное описание способов синхронизации времени устройства приведено в документации \ManualUnitGeneral.
\end{enumerate}

%-----------------------------------------------
% Система самодиагностики
%-----------------------------------------------
\item Система самодиагностики
\begin{enumerate}
\item Система самодиагностики устройства выполняет тесты в полном объеме при включении устройства, а также непрерывно в фоновом режиме.
Системой самодиагностики устройства в автоматическом режиме контролируются следующие параметры:
\begin{itemize}
\item состояние аппаратной части устройства, в том числе: АЦП модулей ввода аналоговых сигналов, блока питания, ОЗУ, ПЗУ, процессора, модулей дискретных входов, модулей дискретного управления, вспомогательных модулей;
\item температурный режим устройства;
\item состояние синхронизации времени;
\item сохранность исполняемого программного кода (целостность ПО);
\item состояние измерительных цепей;
\item исправность используемых каналов связи.
\end{itemize}
\item Система самодиагностики фиксирует критические неисправности («Неисправность») и некритические ошибки («Ошибка»).
По факту выхода из режимов критической неисправности и предупредительной сигнализации производится запись в журнале событий.
Подробное описание системы самодиагностики приводится в документации \ManualUnitGeneral.
\item В случае появления критической неисправности активируется внутренний сигнал «Критическая неисправность», производится запись в журнал событий, замыкаются контакты реле IRF, прекращается работа алгоритмов основных функций в составе устройства, это состояние устройства сигнализируется свечением светодиода «Неисправность» и отсутствием свечения светодиода «Готовность».
Для выхода из режима критической неисправности требуется перезапуск устройства.
\item В случае появления некритической ошибки активируется внутренний сигнал «Некритическая ошибка», производится запись в журнал событий, алгоритмы основных функций в составе устройства продолжают выполняться, это состояние устройства сигнализируется свечением светодиода «Неисправность» и отсутствием свечения светодиода «Готовность».
Устройство может выйти из режима предупредительной сигнализации без активных действий со стороны пользователя по факту пропадания критериев режима.
\end{enumerate}

%-----------------------------------------------
% Аутентификация и авторизация
%-----------------------------------------------
\item Аутентификация и авторизация пользователей
\begin{enumerate}
\item Устройство имеет встроенную реализацию правил разграничения прав доступа, для этого в устройстве существует ролевая модель разграничения прав доступа.
В устройстве предусмотрены следующие роли:
\begin{itemize}
\item Гость (пароль не требуется, доступен просмотр данных устройства);
\item Дежурный;
\item Инженер;
\item Офицер (сотрудник) безопасности;
\item Администратор (пароль «АААА» (кириллица), учетная запись предназначена для первоначального конфигурирования штатного профиля пользователей).
\end{itemize}
\item Все пользователи делятся на непривилегированных (роль Гость) и привилегированных (все остальные роли).
\item Предустановленные пользователи с именами Гость и Администратор не подлежат удалению и принудительной блокировке, их роли не подлежат изменению.
Гость используется как роль по умолчанию для доступа к просмотру данных устройства в режиме чтение без запроса аутентификационных данных пользователя.
Администратор необходим для добавления/удаления других пользователей, установки им соответствующих ролей и установке/изменении паролей пользователей.
\item После перезагрузки устройства подключенная сессия привилегированного пользователя сбрасывается.
Устройство запускается с сеансом непривилегированного пользователя (роль Гость).
Подробное описание ролей и доступных полномочий и ограничений для каждой из ролей приведено в документации \ManualUnitGeneral.
\item Для возможности работы с устройством пользователь должен пройти аутентификацию на устройстве.
Аутентификация проводится по имени и паролю, указанными пользователем.
\item Функции безопасности устройства предъявляют следующие требования к паролям пользователей:
\begin{itemize}
  \item пароль может содержать буквы и/или цифры;
  \item длина пароля не может быть меньше семи символов.
\end{itemize}
\item Хранение паролей осуществляется в памяти устройства в нечитаемом виде, не позволяющем восстановить пароль.
\item Функции безопасности не дают возможность повторного применения ранее использованных паролей.
По умолчанию глубина хранения ранее заданных пользователем паролей в устройстве составляет четыре пароля.
\item В процессе ввода пароля пользователю отображается только вводимый символ, ранее введенные символы не отображаются в открытом виде.
\item Функции безопасности устройства ограничивают количество неуспешных попыток аутентификации за установленный период времени.
Количество попыток можно задать в диапазоне от 1 до 9 (по умолчанию – 3).
\item В случае превышения количества неуспешных попыток аутентификации:
\begin{itemize}
  \item выполняется блокировка доступа пользователя к устройству. Длительность блокировки определяется в настройках устройства пользователем с ролью Администратор в диапазоне от 1 минуты до 9999 минут (по умолчанию – 30 минут);
  \item формируется соответствующее событие журнала событий безопасности;
  \item формируется сигнализация для АСУ ТП о превышении заданного количества неуспешных попыток доступа.
\end{itemize}
\item В случае, когда количество неуспешных попыток аутентификации пользователя больше нуля, но меньше максимально допустимого, оно будет сброшено для данного пользователя по истечении времени сброса неуспешных попыток (5 минут).
\item Блокировка возможности аутентификации пользователя снимается в случае:
\begin{itemize}
  \item истечения времени блокировки входа;
  \item принудительной разблокировки пользователем с ролью Администратор.
\end{itemize}
\item Пользователь может быть заблокирован принудительно пользователем с ролью Администратор. Разблокирование данного пользователя может осуществить так же только пользователь с ролью Администратор.
\item Все настраиваемые параметры безопасности (время неактивности пользователя, количество безуспешных попыток входа, время блокировки входа) разрешены для редактирования пользователям с ролью офицера безопасности. При сбросе устройства к заводским настройкам, параметры безопасности сохраняются.

% Предупреждение о потере пароля администратора
\vspace{-0.3cm}
\noindent
\begin{tabularx}{\linewidth}{p{2em} X}
  \begin{minipage}[c][3\baselineskip][b]{2em}
    \centering\textcolor{unidarkgreen}{\Huge\faCircleInfo}
  \end{minipage} &
  \hyphenpenalty=0 \textbf{При утере пароля учетной записи пользователя с ролью Администратор, восстановление административных прав возможно только путем перепрограммирования устройства специалистом предприятия-изготовителя, что повлечет за собой сброс настроечных параметров на заводские.}
\end{tabularx}
\vskip 2mm

\end{enumerate}

%-----------------------------------------------
% Блок интерфейсный
%-----------------------------------------------
\item\label{sec1:hmi} Блок интерфейсный (ЮНИТ-ИЧМ)
\begin{enumerate}
\item Для местного управления и контроля за состоянием устройства предусматривается подключение блока интерфейсного «ЮНИТ-ИЧМ» производства ООО «Юнител Инжиниринг».
Блок подключается к «ЮНИТ-М500» посредством патч-корда с разъемами RJ45. В устройстве «ЮНИТ-М500» для подключения к блоку может использоваться отдельный порт X5, расположенный на модуле центрального процессора или любой из задействованных портов связи Ethernet (X1, Х2, Х3, X4).
\item Более подробная информация по работе с устройством «ЮНИТ-ИЧМ» приведена в документации \ManualUnitSRZA.
\item Подробное описание блока приведено в документе \ManualUnitHMI. Блок интерфейсный «ЮНИТ-ИЧМ» является опциональным для заказа.
\end{enumerate}

\end{enumerate}

%===============================================================
% 7. Организация обмена информацией
%===============================================================
{\color{unidarkgreen}\subsection{Организация обмена информацией}}

\begin{enumerate}
\item \sloppy Описание и настройки доступных портов и интерфейсов связи приводится в документации \ManualUnitGeneral.
\item В устройствах с поддержкой протокола передачи данных МЭК 61850 в качестве входных сигналов ФСУ могут быть использованы данные из GOOSE-сообщений, на которые оформлена подписка.
Количество сигналов GOOSE, получаемых от интеллектуальных устройств, в данном типоисполнении устройства равно 12 шт.
\item Сигналы из GOOSE-сообщений проходят предварительную обработку перед тем как попасть в «Блок объединительный» и далее в ФСУ (см. приложение \ref{app:fsu}).
В предварительную обработку входят следующие операции:
\begin{itemize}
\item проверка качества принимаемого сигнала;
\item алгоритм обработки резервирования GOOSE-сообщений;
\item алгоритм обработки симуляции GOOSE-сообщений;
\item алгоритм обработки метки «Test» GOOSE-сообщений;
\item подстановка значения по умолчанию при «плохом» качестве данных, принимаемых из GOOSE-сообщений.
\end{itemize}
Подробнее о подписке на GOOSE-сообщения указано в \ManualUnitSRZA.
\item В устройствах с поддержкой протокола передачи данных МЭК 61850 для связи со смежными устройствами может быть использован протокол GOOSE по МЭК 61850-8.1.
Количество сигналов GOOSE-сообщений, передаваемых в смежные интеллектуальные устройства, зависит от типоисполнения ЮНИТ-М500, первичного оборудования и в общем случае составляет не менее 100 шт.
Подробнее о публикации GOOSE-сообщений указано в \ManualUnitSRZA.
\end{enumerate}

%===============================================================
% 8. Прочие подразделы
%===============================================================
{\color{unidarkgreen}\subsection{Средства измерения, инструмент и принадлежности}}
Перечень оборудования и средств измерения, рекомендуемый для проведения эксплуатационных проверок устройства, приведен в документации \ManualUnitGeneral.

{\color{unidarkgreen}\subsection{Маркировка и пломбирование}}
Сведения по маркировке и пломбировании устройства приведены в документации \ManualUnitGeneral.

{\color{unidarkgreen}\subsection{Упаковка}}
Сведения по упаковке устройства приведены в документации \ManualUnitGeneral.